% Supplement to p1.tex. Its theorem, section and equation numbers are those of p1.pdf, read from p1.aux;
% build p1.tex first (make papers does).
\newcommand{\papertitle}{Supplement to ``A log-concavity criterion for finite nonnegative linearization'': formalization and AI use}
\newcommand{\papershort}{SUPPLEMENT: A LOG-CONCAVITY CRITERION FOR FINITE LINEARIZATION}
\newcommand{\paperdate}{November 1, 2026}
\documentclass[twoside,10pt]{article}
\usepackage[T1]{fontenc}
\usepackage{lmodern}
\usepackage{amsmath,amssymb,amsthm,amsfonts,mathtools}
\usepackage[letterpaper,total={5.5in,8in},vmarginratio=1:1,hmarginratio=1:1,headsep=21pt]{geometry}
\usepackage{xcolor}
\usepackage{array,longtable}
\usepackage{xr-hyper}
\usepackage[breaklinks=true,linktocpage=true]{hyperref}
\hypersetup{colorlinks=true,linkcolor=blue!45!black,urlcolor=blue!45!black,citecolor=blue!45!black,
pdfauthor={Leslie P. Polzer},pdftitle={\papertitle},
pdfsubject={Finite orthogonal polynomials and nonnegative linearization},
pdfkeywords={orthogonal polynomials, Krawtchouk polynomials, nonnegative linearization, quadratic transformation, hypergroups, Lean}}
\medmuskip=4mu
\thickmuskip=8mu
\setlength{\emergencystretch}{2em}
\numberwithin{equation}{section}
\newtheoremstyle{slbody}{8pt}{8pt}{\normalfont\slshape}{}{\bfseries}{.}{0.5em}
{\thmname{#1}\thmnumber{ #2}\thmnote{ {\normalfont\itshape(\ignorespaces#3\unskip)}}}
\theoremstyle{slbody}
\newtheorem{theorem}{Theorem}[section]
\newtheorem{lemma}[theorem]{Lemma}
\newtheorem{corollary}[theorem]{Corollary}
\newtheorem{proposition}[theorem]{Proposition}
\theoremstyle{definition}
\newtheorem{definition}[theorem]{Definition}
\theoremstyle{remark}
\newtheorem{remark}[theorem]{Remark}
\renewcommand{\qedsymbol}{$\blacksquare$}
\newcommand{\qint}[1]{[#1]_q}
% \xnum{P1-thm:lc}{3.1} prints the fixed number of a result in a companion paper; check_xref.py checks it.
\newcommand{\xnum}[2]{#2}
\title{\papertitle}
\author{Leslie P. Polzer}
\newcommand{\authoraffiliation}{Independent Researcher}
\newcommand{\authoremail}{polzer@fastmail.com}
\date{\paperdate}
\makeatletter
\renewcommand{\section}{\@startsection{section}{1}{\z@}{-3.5ex \@plus -1ex \@minus -.2ex}{2.3ex \@plus .2ex}{\normalfont\large\bfseries}}
\renewcommand{\@seccntformat}[1]{\csname the#1\endcsname.\enspace}
\def\tagform@#1{\maketag@@@{(\ignorespaces{\oldstylenums{#1}}\unskip\@@italiccorr)}}
\renewcommand{\eqref}[1]{\textup{{\normalfont(\oldstylenums{\ref{#1}}}\normalfont)}}
\newcommand{\shorttitle}{\papershort}
\newcommand{\shortauthors}{LESLIE P. POLZER}
\newcommand{\ps@bookheader}{%
\renewcommand\@oddfoot{\hfil}\renewcommand\@evenfoot{\hfil}%
\renewcommand\@oddhead{\ifnum\value{page}>1{\small\hfil\shortauthors}\hfil\thepage\else\hfil\fi}%
\renewcommand\@evenhead{\thepage\hfil\small\shorttitle\hfil}}
\pagestyle{bookheader}
\renewcommand{\maketitle}{\begin{center}{\large\bfseries\@title}\\\vskip36pt
{\scshape\@author}\par\medskip{\small\authoraffiliation}\par\smallskip
{\small\texttt{\authoremail}}\par\smallskip{\small\@date}\par\vskip36pt\end{center}}
\makeatother
\renewenvironment{abstract}{\quotation\noindent\small{\bfseries\abstractname.\enspace}}{\endquotation}

\externaldocument[M-]{p1}
\begin{document}
\maketitle

\noindent This supplement to~\cite{main} gives the Lean declaration for each of its results, the statements that are not formalized, and the full AI-use disclosure. Theorem, section, equation and appendix numbers refer to~\cite{main}.

\section*{Code and formalization}
The companion code~\cite{code} contains a Lean~4 development and exact-arithmetic diagnostics for every result. Its axiom audit rejects dependencies outside \path{propext}, \path{Classical.choice} and \path{Quot.sound}; all declarations below have compiled proofs without added axioms or admitted goals. The documentation revision~\cite{docs} describes the diagnostics and explains how each Lean statement matches the text.
\begin{longtable}{@{}>{\raggedright\arraybackslash}p{0.36\linewidth}>{\raggedright\arraybackslash}p{0.59\linewidth}@{}}
Result & Lean declaration\\[3pt]
\endhead
Lemma~\ref{M-lem:finite}: finite realization & \path{Szkraw.finite_realization}, \path{Szkraw.product_triangle_zero}\\[3pt]
Theorem~\ref{M-thm:lc}: log-concave Newton criterion & \path{Szkraw.logconcave_criterion}\\[3pt]
Remark~\ref{M-rem:bernstein} & \path{Szkraw.lc_twoParam}, \path{Szkraw.lc_bernstein}, \path{Szkraw.constantRatio_twoParam}\\[3pt]
Corollary~\ref{M-cor:kraw}: Krawtchouk range & \path{Szkraw.full_krawtchouk_binomial_range}, \path{Szkraw.binomial_normalized_at_zero}, \path{Szkraw.binomial_node_probability_convention}, \path{Szkraw.krawtchouk_range_sharp}\\[3pt]
Corollary~\ref{M-cor:additional}: geometric arrays & \path{Szkraw.additional_geometric_class}\\[3pt]
Theorem~\ref{M-thm:decomp}: quadratic decomposition & \path{Szkraw.decomp_positivity}, \path{Szkraw.decomp_converse}, \path{Szkraw.even_block}, \path{Szkraw.odd_block}, \path{Szkraw.evenComp_admissible}, \path{Szkraw.oddComp_admissible}\\[3pt]
Corollary~\ref{M-cor:parity}: parity-split criterion & \path{Szkraw.parity_split}, \path{Szkraw.parity_split_sharp}, \path{Szkraw.ps_even_params}, \path{Szkraw.ps_odd_params}\\[3pt]
Theorem~\ref{M-thm:szwarc}: Szwarc's quadratic system & \path{Szkraw.szwarc_full}, \path{Szkraw.szwarc_sharp}, \path{Szkraw.realization_symmetric}\\[3pt]
Lemma~\ref{M-lem:corner-value}: corner value & \path{Szkraw.corner_value}\\[3pt]
Proposition~\ref{M-prop:raise}: raising the top mass & \path{Szkraw.raise_measure}, \path{Szkraw.raise_corner}, \path{Szkraw.raise_sign}, \path{Szkraw.raise_bound}\\[3pt]
Corollary~\ref{M-cor:radius-exp} and its no-uniform-radius conclusion & \path{Szkraw.radius_exp}, \path{Szkraw.radius_exp_asymp}, \path{Szkraw.radius_no_uniform}\\[3pt]
The two examples in the paragraph ``How sharp is log-concavity?'' and Appendix~\ref{M-app:search} ($\mu=\frac1{10}$ and $\mu=\frac1{100}$) & \path{Szkraw.lc_violation_holds}, \path{Szkraw.lc_violation_fails}\\[3pt]
The combinations $(1-\lambda)K_{N,r}+\lambda G_{N,q,r}$ in the paragraph ``How sharp is log-concavity?'' and the seven failures listed in Appendix~\ref{M-app:search} & \path{Szkraw.combo_shape}, \path{Szkraw.combo_fails_4_5_10}, \path{Szkraw.combo_fails_3_5_20}, \path{Szkraw.combo_fails_1_2_20}, \path{Szkraw.combo_fails_2_5_60}, \path{Szkraw.combo_fails_1_3_100}, \path{Szkraw.combo_fails_3_10_200}, \path{Szkraw.combo_fails_2_7_200}\\[3pt]
Footnote before Corollary~\ref{M-cor:kraw}: the parameter convention & \path{Szkraw.szwarc_footnote}\\[3pt]
Remarks after Theorem~\ref{M-thm:lc}: Newton coefficients of the examples & \path{Szkraw.newton_examples}, \path{Szkraw.convex_counterexamples_newton_signs}\\[3pt]
After Corollary~\ref{M-cor:additional}: Koelink's recurrence, transcribed by hand; agreement checked in Lean & \path{Szkraw.koelink_recurrence_71}\\[3pt]
Section~\ref{M-sec:szwarc}: the hypotheses of Szwarc's Theorem~3.3, transcribed from~\cite{szwarc}, for \eqref{M-eq:szwarc-q} & \path{Szkraw.szwarc_thm33_check}
\end{longtable}
The entry for Theorem~\ref{M-thm:lc} formalizes Lemmas~\ref{M-lem:lc-ineq} and~\ref{M-lem:detours}, the operator induction, and Lemma~\ref{M-lem:kernel-rec} in converse form: the solution of~\eqref{M-eq:kernel-rec} yields the linearization coefficients. As stated in the introduction of~\cite{main}, the Lean proofs of Corollaries~\ref{M-cor:kraw} and~\ref{M-cor:additional} obtain positivity from the geometric-increment criterion of~\cite{P3} rather than from Theorem~\ref{M-thm:lc}, and so does the Lean proof of Corollary~\ref{M-cor:parity} for its two components; Theorem~\ref{M-thm:szwarc} is derived in Lean from Corollary~\ref{M-cor:parity}. The statements are the ones given here. The Lean proof of Lemma~\ref{M-lem:corner-value} uses the algebra of~\cite{P3}. Thus these formal proofs use the construction of~\cite{P3}, whereas the written proofs in this paper deduce the corollaries from Theorem~\ref{M-thm:lc}. The written proof of Theorem~\ref{M-thm:lc} does not use~\cite{P3}, and its Lean proof follows the written route.

Not formalized: the comparisons with the criteria of Askey~\cite{askey}, of Szwarc~\cite{szwarc2003} and of Miclo~\cite{miclo,miclo2019} and of M{\l}otkowski and Szwarc~\cite{mlotszwarc} in the introduction and with the quadratic transformations of Chihara and Stanton~\cite{chiharastanton}, Szwarc~\cite{szwarc1992,szwarc1996} and Kahler~\cite{kahler} in Section~\ref{M-sec:szwarc}, which concern results of the literature; the example after Theorem~\ref{M-thm:lc}, an instance of that theorem whose arithmetic is not checked separately; Koelink's construction and the dual $q$-Krawtchouk interpretation after Corollary~\ref{M-cor:additional}, and the identification with distance-regular graphs that follows it (only his recurrence, transcribed by hand, is compared in Lean, and Lean does not check the transcription; the same applies to Szwarc's Theorem~3.3 in the table); and the float and high-precision searches concerning $V_N$ in Appendix~\ref{M-app:search}, which give numerical evidence over sampled parameters, not certified thresholds. More generally, attributions to other works and statements about what they contain are not formalized, for instance the earlier treatment of the Krawtchouk range by Eagleson and Szwarc: Lean checks our statements, not the cited papers.

\section*{AI-use disclosure}
Three AI models were used: Claude Opus 5.5 (Anthropic), run in the Claude Code agent; GPT-6 Astra or GPT-6.1 Sol (OpenAI; which of the two was not recorded), run in the Codex agent; and GPT-6 Pro (OpenAI), used in ChatGPT. Claude Opus 5.5 and the Codex model were used to develop proof arguments, check source passages, construct exact-arithmetic diagnostics, draft and revise the paper, and prepare the Lean formalization. The kernel recurrence and the operator inequality in the proof of Theorem~\ref{M-thm:lc} arose in research consultations with GPT-6 Pro; the closed form of the recurrence for every degree and the exact checks were completed with Claude Opus 5.5, which also developed Sections~\ref{M-sec:szwarc} and~\ref{M-sec:limits} and prepared the Lean formalization of all numbered results. GPT-6 Pro reviewed the results over several rounds, as part of a longer manuscript and with independent exact checks, and the errors it found were corrected. For Theorem~\ref{M-thm:lc} the reviewer is therefore the model that proposed its key steps. This AI review is not human peer review, and no independent human verification is claimed. The author set the research objective and is responsible for publication decisions and the final mathematical claims. ``Machine-verified'' refers only to the Lean theorems listed in the section on code and formalization, with the axioms stated there; other remarks and computational observations are not claimed to be machine-verified.

\begin{thebibliography}{99}

\bibitem{main} L. P. Polzer, \emph{A log-concavity criterion for finite nonnegative linearization}, 2026, in release~1.0.6 of the companion code~\cite{code}, \url{https://github.com/skypher/szkraw/releases/tag/v1.0.6}.
\bibitem{askey} R. Askey, \emph{Three notes on orthogonal polynomials}, Stichting Mathematisch Centrum, Afdeling Toegepaste Wiskunde, report TW~114/69, Amsterdam, October 1969; Part~II: \emph{Linearization of the product of orthogonal polynomials II}, \url{https://ir.cwi.nl/pub/7741}.
\bibitem{chiharastanton} L. Chihara and D. Stanton, \emph{Association schemes and quadratic transformations for orthogonal polynomials}, Graphs and Combinatorics \textbf{2} (1986), no.~1, 101--112, \href{https://doi.org/10.1007/BF01788084}{\nolinkurl{doi:10.1007/BF01788084}}.
\bibitem{kahler} S. Kahler, \emph{Nonnegative and strictly positive linearization of Jacobi and generalized Chebyshev polynomials}, Constructive Approximation \textbf{54} (2021), no.~2, 207--236, \href{https://doi.org/10.1007/s00365-021-09552-3}{\nolinkurl{doi:10.1007/s00365-021-09552-3}}; preprint \url{https://arxiv.org/abs/1812.05542}.
\bibitem{miclo} L. Miclo, \emph{On the hypergroup property}, Annales de la Facult\'e des Sciences de Toulouse, s\'erie~6, \textbf{26} (2017), no.~2, 417--435, \href{https://doi.org/10.5802/afst.1539}{\nolinkurl{doi:10.5802/afst.1539}}.
\bibitem{miclo2019} L. Miclo, \emph{On the Markov commutator}, Bulletin des Sciences Math\'ematiques \textbf{154} (2019), 1--35, \href{https://doi.org/10.1016/j.bulsci.2019.02.003}{\nolinkurl{doi:10.1016/j.bulsci.2019.02.003}}; preprint arXiv:1505.02540v1 (2015), to whose numbering the theorem locator refers.
\bibitem{mlotszwarc} W. M{\l}otkowski and R. Szwarc, \emph{Nonnegative linearization for polynomials orthogonal with respect to discrete measures}, Constructive Approximation \textbf{17} (2001), no.~3, 413--429, \href{https://doi.org/10.1007/s003650010039}{\nolinkurl{doi:10.1007/s003650010039}}.
\bibitem{P3} L. P. Polzer, \emph{Exact and perturbative criteria for finite nonnegative linearization}, companion paper (2026), in release~1.0.6 of the companion code~\cite{code}, \url{https://github.com/skypher/szkraw/releases/tag/v1.0.6}.
\bibitem{code} L. P. Polzer, \emph{Companion code for finite nonnegative linearization}, release~1.0.6 (2026), \url{https://github.com/skypher/szkraw/releases/tag/v1.0.6}. This release is the single landing point for the evidence: it contains the three papers and their supplements, the Lean development and the scripts, which are unchanged since revision \texttt{2f96f4d}, \url{https://github.com/skypher/szkraw/tree/2f96f4d228215b35e31729f1d1526cd31879f161}, and the documentation of~\cite{docs}.
\bibitem{docs} L. P. Polzer, documentation of the companion code~\cite{code}: \path{papers/FORMALIZATION.md}, \path{papers/DIAGNOSTICS.md} and the record \path{runs/verify/RESULTS.md} of a run of all scripts from a fresh clone of revision \texttt{2f96f4d}, as of revision \texttt{4b2d7e1}, \url{https://github.com/skypher/szkraw/tree/4b2d7e13d5c9a9fe8bbd406959d0ecd2cc192de8}.
\bibitem{szwarc1992} R. Szwarc, \emph{Orthogonal polynomials and a discrete boundary value problem~II}, SIAM Journal on Mathematical Analysis \textbf{23} (1992), no.~4, 965--969, \href{https://doi.org/10.1137/0523053}{\nolinkurl{doi:10.1137/0523053}}.
\bibitem{szwarc1996} R. Szwarc, \emph{Nonnegative linearization and quadratic transformation of Askey--Wilson polynomials}, Canadian Mathematical Bulletin \textbf{39} (1996), no.~2, 241--249, \href{https://doi.org/10.4153/CMB-1996-031-x}{\nolinkurl{doi:10.4153/CMB-1996-031-x}}.
\bibitem{szwarc2003} R. Szwarc, \emph{A necessary and sufficient condition for nonnegative product linearization of orthogonal polynomials}, Constructive Approximation \textbf{19} (2003), no.~4, 565--573, \href{https://doi.org/10.1007/s00365-002-0524-z}{\nolinkurl{doi:10.1007/s00365-002-0524-z}}.
\bibitem{szwarc} R. Szwarc, \emph{Orthogonal polynomials and Banach algebras}, in \emph{Inzell Lectures on Orthogonal Polynomials}, W. zu Castell, F. Filbir, and B. Forster (eds.), Advances in the Theory of Special Functions and Orthogonal Polynomials, vol.~2, Nova Science Publishers, 2005, 103--139; \S5, Problem~2. Author version (2004, 97--133; Problem~2 on printed pp.~125--126): \url{https://www.math.uni.wroc.pl/~szwarc/pdf/inzell.pdf}.
\end{thebibliography}
\end{document}
